% Official source: https://www.imo-official.org/assets/documents/problems/2015/2015_eng.pdf
\competitionyear{2015}

\begin{problem}{2015}{1}
A finite set $S$ of points in the plane is called \emph{balanced} if, for every pair of distinct points $A, B\in S$, some $C\in S$ satisfies $AC=BC$. It is called \emph{centre-free} if no point $P\in S$ satisfies $PA=PB=PC$ for three distinct points $A, B, C\in S$.
\begin{problemsubparts}
    \item Prove that a balanced set of $n$ points exists for every integer $n\ge3$.
    \item Determine all integers $n\ge3$ for which a balanced centre-free set of $n$ points exists.
\end{problemsubparts}
\end{problem}

\begin{problem}{2015}{2}
Find all triples $(a, b, c)$ of positive integers such that $ab-c$, $bc-a$, and $ca-b$ are all powers of $2$.
\end{problem}

\begin{problem}{2015}{3}
Let $ABC$ be an acute triangle with $AB>AC$, circumcircle $\Gamma$, orthocentre $H$, and altitude foot $F$ from $A$. Let $M$ be the midpoint of $BC$. Let $Q$ be the point on $\Gamma$ satisfying $\angle HQA=90^\circ$, and let $K$ be the point on $\Gamma$ satisfying $\angle HKQ=90^\circ$. Assume that $A, B, C, K, Q$ are distinct and appear on $\Gamma$ in that order.

Prove that the circumcircles of $KQH$ and $FKM$ are tangent.
\end{problem}

\begin{problem}{2015}{4}
Let $\Omega$ be the circumcircle of triangle $ABC$, with circumcentre $O$. A circle $\Gamma$ centred at $A$ meets segment $BC$ at distinct points $D, E$ in the order $B, D, E, C$. Let $F, G$ be the intersections of $\Gamma$ with $\Omega$ such that $A, F, B, C, G$ occur in this order on $\Omega$. The circumcircle of $BDF$ meets segment $AB$ again at $K$, and the circumcircle of $CGE$ meets segment $CA$ again at $L$.

Suppose that lines $FK$ and $GL$ are distinct and intersect at $X$. Prove that $X$ lies on $AO$.
\end{problem}

\begin{problem}{2015}{5}
Find all functions $f\colon\R\to\R$ for which
\[
    f(x+f(x+y))+f(xy)=x+f(x+y)+yf(x)
\]
for every $x, y\in\R$.
\end{problem}

\begin{problem}{2015}{6}
An integer sequence $(a_j)_{j\ge1}$ satisfies
\begin{problemclauses}
    \item $1\le a_j\le2015$ for every $j\ge1$;
    \item $k+a_k\ne\ell+a_\ell$ whenever $1\le k<\ell$.
\end{problemclauses}
Prove that there exist positive integers $b, N$ such that
\[
    \left|\sum_{j=m+1}^{n}(a_j-b)\right|\le1007^2
\]
for all integers $n>m\ge N$.
\end{problem}
